Recognizing the limitations of existing platforms and the significant impact of dietary restrictions, we believe our web application offers a compelling solution for individuals with severe allergies or strict dietary needs. By facilitating effortless product substitutions, our platform empowers users to maintain a healthy and balanced diet while enjoying their favorite foods. This is particularly impactful for those with allergies, as it helps them avoid potentially life-threatening health risks. Ultimately, our goal is to make dietary restrictions less restrictive for all users, enabling them to savor their preferred meals without compromise. This driving motivation has fueled our development efforts.
{\par}
Working on this project has provided our team with invaluable experience and knowledge beyond the technical aspects of implementation and technology utilization. We have gained crucial insights into project planning, design, and the vital importance of effective team communication and task division in today's professional environment. Our project journey encompassed three distinct phases: business analysis, implementation and verification, and finally, testing and validation. During the initial business analysis phase, we conducted thorough market research to assess the need for such an application and identify any existing competitors. Upon confirming a market gap, we proceeded to gather user requirements through the creation of user personas and a functional prototype. We also established a product backlog and adopted the Agile Scrum methodology to guide the subsequent development stages. The implementation phase involved selecting a suitable development methodology, building the necessary database, and embarking on a series of two-week sprints. Following the successful completion of the implementation phase, rigorous testing was conducted to ensure the application's functionality and user experience. These comprehensive testing efforts concluded the project with satisfactory results.
{\par}
Our project could also offer other functionalities to consider for application development. For example, we could enhance user functionality by adding features such as the ability to flag disliked ingredients. Another valuable addition would be allergy-based filtering, allowing users to automatically swap allergens in recipes. Additionally, the implementation of a recommendation system could significantly improve the user experience by suggesting recipes based on individual preferences. Advanced filtering methods could also be introduced that could dynamically adjust ingredients in recipes to match the desired diet. Finally, the inclusion of an admin interface and dynamic streams would provide better system management and real-time updates. These improvements would not only enhance the user experience, but also improve the administration of the application.
    